\section{Discussion}
\label{sec:discussion}

A finite collection of game instances can be certified by exhaustive
computation; a general knot-theoretic claim cannot. That distinction structures
everything above: \texttt{SOLVED: YES} certifies a finite minimax table, while
a reusable tactical motif---say, a local crossing pattern whose resolution is
provably optimal across a shadow family, or an MCTS principal variation that
recurs across knot types---would require a mathematical argument or a formally
justified reduction before it becomes a lemma. The Lean formalization of PD
signs, flips, and the crossing-resolution game is the scaffold for such
reductions; the solved trefoil, figure-eight, and $5_2$ games it already
certifies are the base cases.

The trained agent's principal variations are the concrete hunting ground: every
optimal line the exact tables endorse is a candidate motif awaiting
generalization. The diffusion operators of Section~\ref{sec:representations}
offer a complementary, global view---strategic similarity at scale---but their
pilot organization by crossing count may reflect diagram size as much as
strategy, which is precisely what the control checklist is designed to
disentangle.

\subsection{Limitations}
\label{sec:limitations}

\begin{itemize}
    \item Supervised capacity only; no self-play solution yet. The AlphaZero
        run reaches 71.73\% optimal top-policy play on the fixed shadow.
    \item One table diagram per knot type; knot-level claims need multi-diagram
        robustness and shadow-level dedup beyond the current manifest.
    \item Eight-crossing production bound not yet raised (independent
        validation pending); larger crossing numbers need validated terminal
        classification and new exact-table budgets ($3^9 = 19{,}683$ states per
        shadow at nine crossings).
    \item Diffusion geometry exploratory; all quantitative controls pending.
    \item Two of three shared-capacity seeds measured (solving epochs 190 and
        230); the third is queued.
\end{itemize}

\subsection{Next steps}
\label{sec:next}

In near-term order (ROADMAP): multi-seed mixed-size capacity with variance
reporting; independent eight-crossing invariant validation and bound decision;
richer diagram/shadow metadata with additional non-tabulated diagrams under
frozen splits; a shared variable-size AlphaZero self-play run graded by raw
versus MCTS policy against exact tables; canonical activation/MCTS export at
all five stages; and diffusion comparison with the full control battery. Each
step lands in \texttt{docs/EXPERIMENTS.md} with commit, job ID, configuration,
elapsed time, exact metrics, and verdict before it enters this paper.
